\documentclass[FM,Proj]{tulthesis}
% tento dokument používá balíky specifické pro XeLaTeX a lze jej přeložit
% jen XeLaTeXem, nemáte-li instalována použitá (komerční) písma, změňte
% nebo vymažte příkazy \set...font na následujících řádcích

% Autor šablony: Pavel Satrapa: http://www.nti.tul.cz/~satrapa/vyuka/latex-tul/

% Autor komentářů, jejich překladů do EN, nastavení BibLaTeXu a aplikace ČSN ISO 690: Jan Koprnický
% http://www.fm.tul.cz/personal/jan.koprnicky

% ENGLISH EXPLANATION
% \documentclass[FM,Dis,EN,fonts,bw,twoside]{tulthesis} % black and white typing, dissertation thesis at FM, written in English with using of TUL Mono font, text is prepared to print on two sides of the paper
% this document uses packages specific for XeLaTeX and it is possible to 
% compile it by XeLaTeX only, if you haven't installed used (commercial) fonts
% change them or erase commands \set...font in following rows
% settings: FM (faculty: FS, FT, FP, EF, FA, FM, FZS a CXI), Dis (type of thesis: BP, DP, Teze, Autoref, Hab, SP, Proj), EN (written in English language), fonts (activation of TUL fonts), bw (black and white)

% Autor of the template tulthesis: Pavel Satrapa: http://www.nti.tul.cz/~satrapa/vyuka/latex-tul/

% Autor of several comments and their translation into English, BibLaTeX settings and CSN ISO 960 citation standard setting: Jan Koprnický
% http://www.fm.tul.cz/personal/jan.koprnicky

% poslední změna / last modification viz/see 00README-JK.txt

\newcommand{\verze}{2.3}

\usepackage{polyglossia}
\setdefaultlanguage{czech} % comment when English is preferred
%\setdefaultlanguage{english} % comment when Czech is preferred

\usepackage{float}
\usepackage{makeidx}
\makeindex
\usepackage{dirtree}
\usepackage{tabularx}

\usepackage{xunicode}
\usepackage{xltxtra}



% příkazy specifické pro tento dokument / specific commands for this document
\newcommand{\argument}[1]{{\ttfamily\color{\tulcolor}#1}}
\newcommand{\argumentindex}[1]{\argument{#1}\index{#1}}
\newcommand{\prostredi}[1]{\argumentindex{#1}}
\newcommand{\prikazneindex}[1]{\argument{\textbackslash #1}}
\newcommand{\prikaz}[1]{\prikazneindex{#1}\index{#1@\textbackslash #1}}
\newenvironment{myquote}{\begin{list}{}{\setlength\leftmargin\parindent}\item[]}{\end{list}}
\newenvironment{listing}{\begin{myquote}\color{\tulcolor}}{\end{myquote}}
\sloppy

% deklarace pro titulní stránku / title page declaration
\TULtitle{ScanWatcher}{ScanWatcher}
\TULauthor{Maksym Rudniev}

% pro bakalářské, diplomové a disertační práce / for bachelor, master theses and dissertation
\TULprogramme{B0613A140005}{Informační technologie}{Information Technology}
\TULbranch{B0613A140005AI}{Aplikovaná informatika}{Applied informatics}
%\TULbranch{1802T008}{Nějaký jiný obor}{Some other branch}
\TULsupervisor{doc. Ing. Josef Chaloupka, Ph.D.}
%\TULconsultant{doc. RNDr. Pavel Satrapa, Ph.D.}
%\TULconsultant{doc. RNDr. Druhý Konzultant, Ph.D.}
%\TULconsultant{doc. RNDr. Třetí Konzultant, Ph.D.}
\TULyear{2025}

% pro habilitační práce / habilitation thesis
%\TULbranch{}{Technická kybernetika}{Technical cybernetics}
%\TULyear{2022}

% Použití bibLateXu, pracuje s ISO stylem
% BibLaTeX settings, works with ISO style
\usepackage[ 
backend=biber
% ,style=iso-authoryear % styl vyžaduje FZS TUL , místo příkazu \cite{} je potřeba využít \parencite{} (sazba kulatých závorek) / style required by FZS TUL use \parencite{} instead of \cite{}
,style=iso-numeric
%,style=numeric
%,sortlocale=cs_CZ
,sorting=nyt %seřazeno jméno, rok, titul
,autolang=other
,bibencoding=UTF8
%,urldate=edtf
,maxcitenames=2 %maximum v textu citovaných jmen
,maxbibnames=3 %maximum v seznamu vyjmenovaných autorů
]{biblatex}
\addbibresource{refTULTHESIS.bib}% vložení seznamu literárních zdrojů v bib formátu / input of references in bib format

% Úprava iso-numeric.bbx v souladu s požadavky TUL hranaté závorky v číslovaném seznamu / Modification of iso-numeric.bbx in accordance with TUL requirements of square brackets in a numbered list
\DeclareFieldFormat{labelnumberwidth}{\mkbibbrackets{#1}}

% Formátování podle pokynů FZS, při využití stylu iso-authoryear, čárka mezi jmény a poslední jméno se spojkou a / special requirements of FZS TUL 
\DeclareDelimFormat{multinamedelim}{\addcomma\space}

\DeclareDelimFormat{finalnamedelim}{%
	\ifnumgreater{\value{liststop}}{2}{\finalandcomma}{}%
	\addspace\bibstring{and}\space}

\DeclareNameAlias{author}{family-given/given-family} 
%%%%%%%%%%%%%%%%%%%%%%%%%%

\usepackage{csquotes} %užití biblatexu hlasí warnings, důvodem může být použití českých uvozovek v citacích! / solving of problems with Czech quotations
\urlstyle{same} %sazba url odkazů stejným fontem jako ostatní text, řešení problémů v zalamování hypertextových odkazů v citacích / url in references setting into the same form as text 

\usepackage{listings} %package for sourcode typesetting
%řešení problémů češtiny např. při sazbě zdrojových kódů / czech special characters in sourcecode typesetting
\makeatletter
\lst@InputCatcodes
\def\lst@DefEC{%
	\lst@CCECUse \lst@ProcessLetter
	^^80^^81^^82^^83^^84^^85^^86^^87^^88^^89^^8a^^8b^^8c^^8d^^8e^^8f%
	^^90^^91^^92^^93^^94^^95^^96^^97^^98^^99^^9a^^9b^^9c^^9d^^9e^^9f%
	^^a0^^a1^^a2^^a3^^a4^^a5^^a6^^a7^^a8^^a9^^aa^^ab^^ac^^ad^^ae^^af%
	^^b0^^b1^^b2^^b3^^b4^^b5^^b6^^b7^^b8^^b9^^ba^^bb^^bc^^bd^^be^^bf%
	^^c0^^c1^^c2^^c3^^c4^^c5^^c6^^c7^^c8^^c9^^ca^^cb^^cc^^cd^^ce^^cf%
	^^d0^^d1^^d2^^d3^^d4^^d5^^d6^^d7^^d8^^d9^^da^^db^^dc^^dd^^de^^df%
	^^e0^^e1^^e2^^e3^^e4^^e5^^e6^^e7^^e8^^e9^^ea^^eb^^ec^^ed^^ee^^ef%
	^^f0^^f1^^f2^^f3^^f4^^f5^^f6^^f7^^f8^^f9^^fa^^fb^^fc^^fd^^fe^^ff%
	% české znaky
	^^^^010c^^^^010d^^^^010e^^^^010f^^^^011a^^^^011b^^^^0147^^^^0148%
	^^^^0158^^^^0159^^^^0160^^^^0161^^^^0164^^^^0165%
	^^^^016e^^^^016f^^^^017d^^^^017e%
	^^00}
\lst@RestoreCatcodes
\makeatother

\begin{document}
	
	\ThesisStart{male}
	%\ThesisStart{zadani-a-prohlaseni.pdf}
	
	\begin{abstractCZ}
	Semestrální projekt popisuje desktopovou aplikaci určenou k pohodlnému prohlížení
	a editaci naskenovaných historických dokumentů. Aplikace umožňuje zobrazení
	naskenovaných stránek spolu s jejich textovým přepisem, přepínání mezi
	českým a německým jazykem, úpravu textu a vyhledávání v obsahu dokumentu.
	Cílem aplikace je usnadnit práci s digitalizovanými knihami a archivy.
	\end{abstractCZ}
	
	\begin{keywordsCZ}
		naskenované dokumenty, prohlížeč, editor textu, digitalizace, archiv
	\end{keywordsCZ}
	
	\vspace{2cm}
	
	\begin{abstractEN}
	This semester project describes a desktop application designed for convenient viewing
	and editing of scanned historical documents. The application allows users
	to display scanned pages alongside their text transcriptions, switch between
	Czech and German language versions, edit the transcribed text, and search
	through the document contents. The goal of the application is to simplify
	working with digitized books and archives.
	\end{abstractEN}
	
	\begin{keywordsEN}
		scanned documents, viewer, text editor, digitization, archive
	\end{keywordsEN}
	
	\clearpage
	
	\begin{acknowledgement}
		Rád bych poděkoval všem, kteří přispěli ke vzniku tohoto dílka.
	\end{acknowledgement}
	
	\tableofcontents
	
	\clearpage
	
	\begin{abbrList}
		\textbf{TUL} & Technická univerzita v~Liberci \\
		\textbf{FM} & Fakulta mechatroniky, informatiky a mezioborových studií
		Technické univerzity v~Liberci\\
		\textbf{FUA} & Fakulta umění a architektury \\
		\textbf{FS} & Fakulta strojní \\
		\textbf{FT} & Fakulta textilní \\
		\textbf{FP} & Fakulta přírodovědně-humanitní a pedagogická \\
		\textbf{EF} & Ekonomická fakulta \\
		\textbf{FZS} & Fakulta zdravotnických studií \\
		\textbf{CXI} & Ústav pro nanomateriály, pokročilé technologie a inovace \\
		\textbf{AI} & Artificial Intelligence 
	\end{abbrList}
	
	
\chapter{Úvod}

Digitalizace historických dokumentů představuje jeden z~klíčových nástrojů
pro zachování kulturního dědictví. Samotné naskenování dokumentu však nestačí
-- pro plnohodnotné vědecké využití je nezbytné, aby byl obsah přístupný
také ve formě editovatelného a prohledávatelného textu.\\[6pt]

Práce s~digitalizovanými dokumenty přináší řadu specifických výzev. Historické
texty jsou často psány v~jazycích, které se od současné normy výrazně liší,
a jejich přepis vyžaduje odborné znalosti i pečlivou manuální kontrolu.
Existující obecné nástroje přitom nejsou pro tento typ práce optimalizovány
-- neumožňují současné zobrazení naskenované předlohy a textového přepisu
ani strukturovanou navigaci v~rozsáhlých dokumentech.\\[6pt]

Cílem tohoto semestrálního projektu je návrh a implementace desktopové aplikace, která
poskytuje specializované prostředí pro prohlížení a editaci naskenovaných
historických dokumentů. Aplikace umožňuje zobrazit naskenovanou stránku
vedle jejího textového přepisu, přepínat mezi českou a německou jazykovou
verzí, editovat text a procházet obsah dokumentu podle ročníků a kapitol.\\[6pt]

V rámci projektu jsou využívána reálná historická data ve formě naskenovaných knih německých jizerských autorů z~období let 1891 až 1936. Tyto materiály představují cenný zdroj informací o regionální historii, kultuře a jazyce, avšak jejich zpracování je komplikováno dobovou typografií, jazykovými odlišnostmi a kvalitou dochovaných skenů. Právě tyto charakteristiky činí zvolená data vhodným testovacím prostředím pro navrhovanou aplikaci a zároveň zdůrazňují potřebu specializovaného nástroje pro efektivní práci s historickými dokumenty.
		
\chapter{Uživatelská dokumentace}
\section{Přehled aplikace}



\begin{figure}[H]
	\centering
	\includegraphics[width=\textwidth-3cm]{images/app_overview}
	\caption{Celkový pohled na aplikaci}
	\label{fig:app_overview}
\end{figure}


Na obrázku \ref{fig:app_overview} je zobrazen \textbf{celkový pohled na aplikaci} v jejím hlavním okně. Rozhraní je rozděleno do několika částí, které umožňují efektivní práci s digitalizovanými dokumenty. V horní části se nachází \textbf{lišta nabídek} obsahující záložky \textbf{File}, \textbf{Edit} a \textbf{Settings}, které poskytují přístup ke všem hlavním funkcím aplikace. Levou část okna tvoří \textbf{Image Viewer}, ve kterém je zobrazena naskenovaná stránka dokumentu. Tento panel umožňuje práci s obrazem, například jeho přiblížení, oddálení nebo posun, a také \textbf{přepínání mezi jazykovými verzemi} předlohy. Pravou část okna zabírá \textbf{textový editor}, který slouží k prohlížení a úpravám přepisu dokumentu. Editor podporuje \textbf{přepínání mezi českou a německou verzí textu}, mezi původním a upraveným obsahem a umožňuje aktivaci \textbf{režimu editace}.


\begin{figure}[H]
	\centering
	\includegraphics[width=\textwidth-1cm]{images/app_diagram}
	\caption{Diagram aplikace}
	\label{fig:app_diagram}
\end{figure}

Na obrázku \ref{fig:app_diagram} je znázorněn \textbf{diagram aplikace}, který popisuje strukturu jednotlivých komponent a jejich vzájemné propojení. Diagram ukazuje, jak jsou hlavní části aplikace --- tedy \textbf{uživatelské rozhraní}, \textbf{modul pro zobrazování obrazů} a \textbf{textový editor} --- organizovány a jak mezi sebou komunikují. Z diagramu je patrné, že aplikace je navržena \textbf{modulárně}, což umožňuje oddělení jednotlivých funkcí a jejich snadnější rozšiřitelnost. Součástí návrhu je také práce s obsahem dokumentu, který je přístupný prostřednictvím okna \textbf{Contents}, a konfigurace aplikace, například \textbf{nastavení písma}, dostupná v sekci \textbf{Settings}.

\clearpage

\section{Main Window}

\subsection{Menu Bar}
Lišta nabídek obsahuje tři záložky: \textbf{File}, \textbf{Edit} a \textbf{Settings}.

\begin{figure}[H]
	\centering
	\begin{minipage}[t]{0.3\textwidth}
		\centering
		\includegraphics[height=3cm]{images/menu_file}
		\caption{Záložka File}
		\label{fig:menu_file}
	\end{minipage}
	\hfill
	\begin{minipage}[t]{0.3\textwidth}
		\centering
		\includegraphics[height=3cm]{images/menu_edit}
		\caption{Záložka Edit}
		\label{fig:menu_edit}
	\end{minipage}
	\hfill
	\begin{minipage}[t]{0.3\textwidth}
		\centering
		\includegraphics[height=3cm]{images/menu_settings}
		\caption{Záložka Settings}
		\label{fig:menu_settings}
	\end{minipage}
\end{figure}

\subsubsection{File}
Záložka \textbf{File} slouží k práci se soubory a umožňuje otevření okna \textbf{Contents} (viz sekci~\ref{sec:contents}).

\begin{description}
	\item[Open Image] Umožňuje testovací prohlížení obrázků
	\item[Open Text] Umožňuje testovací prohlížení textu
	\item[Contents] Otevře okno \textbf{Contents} (viz sekci~\ref{sec:contents})
\end{description}

\subsubsection{Edit}
Záložka \textbf{Edit} obsahuje nástroje pro práci s textem. Většina z nich je dostupná také přímo v rozhraní (viz sekci~\ref{sec:text_editor}).
\label{sec:menu_edit}

\begin{description}
	\item[Start Editing / Stop Editing] Přepínač pro zapnutí nebo vypnutí režimu úprav
	\item[Translate / Stop Translate] Přepínač pro zobrazení nebo skrytí překladu textu
	\item[Save Changes] Uloží provedené změny
	\item[Discard Changes] Zruší provedené změny
	\item[View Original / Exit Original View] Přepínač pro zobrazení nebo skrytí původní verze textu
\end{description}

\subsubsection{Settings}
Záložka \textbf{Settings} v současné době obsahuje pouze konfiguraci písma (viz sekci~\ref{sec:font_settings}).

\begin{description}
	\item[Open Settings] Otevře okno \textbf{Settings} (viz sekci~\ref{sec:settings})
\end{description}

\subsection{Image Viewer}
\label{sec:image_viewer}

\begin{figure}[H]
	\centering
	\includegraphics[height=10cm]{images/image_viewer}
	\caption{Image Viewer}
	\label{fig:image_viewer_main}
\end{figure}

Zobrazuje naskenovaný obrázek stránky dokumentu. \\
Výchozím režimem okna je prohlížeč obrázků. Kromě něj jsou k dispozici dva další režimy zobrazení.

\subsubsection{Image Viewer}

\begin{description}
	\item[Kolečko myši] Přiblížení nebo oddálení obrázku
	\item[Držení LMB + tažení] Posun obrázku
	\item[Posuvník vpravo] Svislý posun obrázku
	\item[Posuvník dole] Vodorovný posun obrázku
\end{description}

\subsubsection{Original German Text a Original Czech Text}

\begin{figure}[H]
	\centering
	\includegraphics[height=10cm]{images/image_viewer_mode}
	\caption{Image Viewer Mode}
	\label{fig:image_viewer_mode}
\end{figure}

Režim \textbf{Original German Text} zobrazuje přepis původního německého textu.
Režim \textbf{Original Czech Text} zobrazuje přepis původního českého textu.

\subsection{Text Editor}
\label{sec:text_editor}

\begin{figure}[H]
	\centering
	\includegraphics[height=10cm]{images/text_editor}
	\caption{Text Editor}
	\label{fig:text_editor}
\end{figure}

Jedná se o vizuální editor textu. V levém dolním rohu se nacházejí tři přepínací tlačítka. Stejné funkce jsou dostupné také v záložce \textbf{Edit} (viz sekci~\ref{sec:menu_edit}).

\begin{description}
	\item[Language (Czech / German)] Přepíná jazyk zobrazeného a editovaného textu
	\item[Version (Edited / Original)] Přepíná mezi upravenou a původní verzí textu
	\item[Mode (Start Editing / Save)] Zapíná režim editace a umožňuje uložení změn
\end{description}

Všechna tři tlačítka lze kombinovat, například je možné editovat německý text.
V režimu \textbf{Original} není editace dostupná.

\section{Contents}
\label{sec:contents}

\begin{figure}[H]
	\centering
	\begin{minipage}[t]{0.48\textwidth}
		\centering
		\includegraphics[height=7cm]{images/contents_window}
		\caption{Okno Contents}
		\label{fig:contents}
	\end{minipage}
	\hfill
	\begin{minipage}[t]{0.48\textwidth}
		\centering
		\includegraphics[height=7cm]{images/contents_search}
		\caption{Vyhledávání v okně Contents}
		\label{fig:contents_search}
	\end{minipage}
\end{figure}

Okno \textbf{Contents} obsahuje tři panely.

\begin{description}
	\item[Search] Textové pole pro vyhledávání podle názvu kapitoly
	\item[Years] Seznam dostupných ročníků
	\item[Chapters] Seznam kapitol vybraného ročníku
\end{description}

Kliknutím na ročník se v panelu Chapters zobrazí odpovídající kapitoly.
Vyhledávání filtruje kapitoly napříč všemi ročníky.

\section{Settings}
\label{sec:settings}

\begin{figure}[H]
	\centering
	\includegraphics[height=10cm]{images/settings_window}
	\caption{Okno Settings}
	\label{fig:settings}
\end{figure}

Okno \textbf{Settings} obsahuje dva panely a dvě tlačítka.

\begin{description}
	\item[Settings list] Výběr kategorie nastavení (v současnosti pouze Font Settings)
	\item[Settings panel] Zobrazuje možnosti vybrané kategorie
	\item[OK] Potvrdí a uloží změny
	\item[Cancel] Zavře okno bez uložení změn
\end{description}

\subsection{Font Settings}
\label{sec:font_settings}

Nastavení písma pro zobrazení textu v Text Editoru (viz sekci~\ref{sec:text_editor}).

\begin{description}
	\item[Font list] Výběr rodiny písma
	\item[Style] Tři zaškrtávací políčka: Bold, Italic, Underline
	\item[Size list] Výběr velikosti písma
	\item[Sample] Náhled textu „Sample“ zobrazený ve zvoleném stylu
\end{description}

\chapter{Technická dokumentace}

\section{Project Structure}

\dirtree{%
	.1 ScanWatcher/.
	.2 app/.
	.2 resources/.
	.3 data/.
	.3 icons/.
	.3 styles/.
	.2 .gitignore.
	.2 conf.json.
	.2 requirements.txt.
}

\subsection{app/}

Složka obsahuje veškerý zdrojový kód projektu.

\begin{table}[H]
	\centering
	\caption{Soubory složky app/}
	\begin{tabularx}{\textwidth}{|l|X|}
		\hline
		\textbf{Soubor} & \textbf{Popis} \\ \hline
		main.py              & Vstupní bod aplikace \\ \hline
		ui\_mainwindow.py    & Definice hlavního okna \\ \hline
		custom\_titlebar.py  & Vlastní titulní lišta (návrh, v aplikaci se nepoužívá) \\ \hline
		page\_controller.py  & Řízení načítání a přepínání stránek \\ \hline
		edit\_controller.py  & Logika režimu úprav \\ \hline
		event\_controller.py & Zpracování událostí uživatelského rozhraní \\ \hline
		image\_viewer.py     & Komponenta prohlížeče obrázků \\ \hline
		contents.py          & Okno obsahu \\ \hline
		contents\_loader.py  & Načítání dat obsahu \\ \hline
		contents\_search.py  & Vyhledávání v obsahu \\ \hline
		settings.py          & Logika okna nastavení \\ \hline
		font\_settings.py    & Nastavení písma \\ \hline
	\end{tabularx}
	\label{tab:app-files}
\end{table}

\clearpage
\subsection{resources/}

Složka obsahuje veškeré prostředky projektu.

\subsubsection{resources/data/}
Hlavní datová složka. Obsahuje dvojici složek pro každý ročník od 1891 do 1941.

\begin{table}[H]
	\centering
	\caption{Soubory datových složek}
	\begin{tabularx}{\textwidth}{|l|l|X|}
		\hline
		\textbf{Složka} & \textbf{Soubor} & \textbf{Popis} \\ \hline
		\texttt{rYYYY/} & \texttt{rYYYY\_sNNN.jpg} & Sken stránky knihy \\ \hline
		\texttt{rYYYY/} & \texttt{rYYYY\_sNNN.txt} & Přepis německého textu \\ \hline
		\texttt{rYYYY/} & \texttt{rYYYY\_sNNN\_new.txt} & Upravená verze přepisu (pokud existuje) \\ \hline
		\texttt{rYYYY\_CZ/} & \texttt{rYYYY\_sNNN\_CZ.txt} & Český překlad německého přepisu \\ \hline
		\texttt{rYYYY\_CZ/} & \texttt{rYYYY\_sNNN\_CZ\_new.txt} & Upravená verze překladu (pokud existuje) \\ \hline
	\end{tabularx}
	\label{tab:data-files}
\end{table}

kde \texttt{YYYY} = rok (4 cifry), \texttt{NNN} = číslo stránky (3 cifry).

\subsubsection{resources/Obsah\_Inhalt\_1891-1941.xls}

Hlavní datový soubor obsahu. Obsahuje seznam kapitol ve formátu uvedeném
v tabulce \ref{tab:xls-columns}.

\begin{table}[H]
	\centering
	\caption{Sloupce souboru Obsah\_Inhalt\_1891-1941.xls}
	\begin{tabularx}{\textwidth}{|l|X|}
		\hline
		\textbf{Sloupec} & \textbf{Popis} \\ \hline
		Rok      & Ročník vydání (1891–1941) \\ \hline
		Stránka  & Číslo stránky kapitoly v daném ročníku \\ \hline
		Kapitola & Název kapitoly \\ \hline
	\end{tabularx}
	\label{tab:xls-columns}
\end{table}

Soubor je načítán při startu okna Contents (viz sekci~\ref{sec:contents})
přes \texttt{contents\_loader.py}.

\subsubsection{resources/icons/}

Složka obsahuje veškeré ikony používané v aplikaci.

\subsubsection{resources/styles/}

Složka obsahuje soubory stylů aplikace.

\subsubsection{main\_dark.qss}
Hlavní tmavý styl aplikace.

\subsection{Kořenové soubory}

\begin{table}[H]
	\centering
	\caption{Kořenové soubory projektu}
	\begin{tabularx}{\textwidth}{|l|X|}
		\hline
		\textbf{Soubor} & \textbf{Popis} \\ \hline
		conf.json        & Konfigurační soubor aplikace (v současné době ukládá pouze nastavení písma) \\ \hline
		requirements.txt & Seznam závislostí Pythonu \\ \hline
		.gitignore       & Pravidla pro ignorování souborů v Gitu \\ \hline
	\end{tabularx}
	\label{tab:root-files}
\end{table}

\clearpage
\section{Scripts}

\subsection{General Scripts}

\subsubsection{app/main.py}

Vstupní bod aplikace. Inicializuje Qt aplikaci, načte styl \texttt{main\_dark.qss},
nastaví ikonu okna a spustí hlavní okno.

\textbf{\texttt{MainController}} — hlavní třída okna, spojuje \texttt{Ui\_MainWindow}
s \texttt{EventController}.

\begin{table}[H]
	\centering
	\caption{Metody třídy MainController}
	\begin{tabularx}{\textwidth}{|l|X|}
		\hline
		\textbf{Metoda} & \textbf{Popis} \\ \hline
		\_\_init\_\_ & Inicializuje UI, vytvoří EventController, nastaví ikonu okna a přepíše closeEvent \\ \hline
	\end{tabularx}
	\label{tab:maincontroller}
\end{table}


\subsubsection{app/event\_controller.py}

Centrální řadič událostí. Propojuje UI signály s logikou aplikace, spravuje
stav načtené stránky a koordinuje ostatní řadiče.


\begin{table}[H]
	\centering
	\caption{Metody třídy EventController}
	\begin{tabularx}{\textwidth}{|l|X|}
		\hline
		\textbf{Metoda} & \textbf{Popis} \\ \hline
		\_\_init\_\_                     & Inicializuje stav, připojí UI signály a načte konfiguraci písma \\ \hline
		\_confirm\_unsaved\_changes       & Zobrazí dialog při neuložených změnách \\ \hline
		\_get\_year\_path                & Vrátí cestu ke složce daného ročníku \\ \hline
		get\_max\_page                   & Zjistí počet stránek ročníku podle .jpg souborů \\ \hline
		\_on\_page\_changed              & Reaguje na změnu stránky a zavolá load\_page \\ \hline
		load\_page                       & Načte obrázek a text pro daný rok, stránku a režim \\ \hline
		openImage                        & Otevře dialog pro výběr obrázku \\ \hline
		openText                         & Otevře dialog pro výběr textového souboru \\ \hline
		openContents                     & Otevře nebo přenese do popředí okno Contents \\ \hline
		onContentsClosed                 & Vymaže referenci na okno obsahu po jeho zavření \\ \hline
		closeEvent                       & Ověří neuložené změny, zavře podokna a ukončí aplikaci \\ \hline
		open\_settings                   & Otevře nebo přenese do popředí okno Settings \\ \hline
		on\_config\_saved               & Aplikuje uložené nastavení písma na editor i prohlížeč \\ \hline
		apply\_font\_config\_at\_startup & Načte conf.json při startu a aplikuje uložené písmo \\ \hline
		highlight\_text                  & Dočasně zvýrazní hledaný text žlutě po dobu 10 sekund \\ \hline
	\end{tabularx}
	\label{tab:eventcontroller}
\end{table}

\clearpage

\subsection{Main Window}

\subsubsection{app/ui\_mainwindow.py}

Definice a sestavení hlavního okna aplikace. Obsahuje dvě třídy.

\textbf{\texttt{EditButtonTextEdit}} — rozšířený \texttt{QTextEdit} se třemi tlačítky
v levém dolním rohu.

\begin{table}[H]
	\centering
	\caption{Metody třídy EditButtonTextEdit}
	\begin{tabularx}{\textwidth}{|l|X|}
		\hline
		\textbf{Metoda} & \textbf{Popis} \\ \hline
		\_\_init\_\_   & Vytvoří tlačítka editButton (tužka), viewButton (oko) a translateButton (DE/CZ) \\ \hline
		resizeEvent    & Při změně velikosti přepočítá pozici tlačítek \\ \hline
	\end{tabularx}
	\label{tab:editbuttontextedit}
\end{table}

\textbf{\texttt{Ui\_MainWindow}} — sestavuje layout hlavního okna a propojuje UI akce.

\begin{table}[H]
	\centering
	\caption{Metody třídy Ui\_MainWindow}
	\begin{tabularx}{\textwidth}{|l|X|}
		\hline
		\textbf{Metoda} & \textbf{Popis} \\ \hline
		setupUi             & Vytvoří celý layout: splitter, PageController a menu bar \\ \hline
		on\_splitter\_moved & Hlídá minimální šířky panelů při přetahování splitteru \\ \hline
	\end{tabularx}
	\label{tab:uimainwindow}
\end{table}

\subsubsection{app/image\_viewer.py}

Komponenta pro zobrazení obrázků a textů stránky. Dědí z \texttt{QGraphicsView}.


\begin{table}[H]
	\centering
	\caption{Metody třídy ImageViewer}
	\begin{tabularx}{\textwidth}{|l|X|}
		\hline
		\textbf{Metoda} & \textbf{Popis} \\ \hline
		\_\_init\_\_              & Inicializuje scénu, scrollbary, QComboBox a QTextBrowser \\ \hline
		get\_event\_controller    & Najde EventController procházením rodičovského stromu \\ \hline
		\_on\_mode\_changed        & Přepne zobrazení podle zvoleného režimu v comboboxu \\ \hline
		\_show\_image\_mode        & Zobrazí obrázek, skryje textový prohlížeč \\ \hline
		\_show\_text\_file         & Načte textový soubor a zobrazí ho v QTextBrowser \\ \hline
		apply\_font\_from\_config  & Aplikuje nastavení písma na QTextBrowser \\ \hline
		load\_image               & Načte .jpg do scény a přizpůsobí měřítko viewportu \\ \hline
		wheelEvent                & Kolečko myši přibližuje / oddaluje obrázek \\ \hline
		resizeEvent               & Přepočítá pozici comboboxu a měřítko obrázku \\ \hline
	\end{tabularx}
	\label{tab:imageviewer}
\end{table}

\clearpage
\subsubsection{app/edit\_controller.py}

Správce stavu textového editoru. Řídí režimy úprav, překladu a zobrazení originálu.


\begin{table}[H]
	\centering
	\caption{Metody třídy TextEditorManager}
	\begin{tabularx}{\textwidth}{|l|X|}
		\hline
		\textbf{Metoda} & \textbf{Popis} \\ \hline
		set\_actions\_enabled     & Zapne nebo vypne všechna tlačítka a akce editoru \\ \hline
		start\_editing            & Přepne editor do režimu úprav, změní ikonu na disketu \\ \hline
		stop\_editing             & Ukončí režim úprav, vrátí ikonu tužky \\ \hline
		toggle\_editing           & Přepínač mezi start\_editing a stop\_editing \\ \hline
		\_confirm\_save\_original & Zobrazí varování při pokusu o uložení přes originál \\ \hline
		save\_changes             & Uloží text do souboru \_new.txt nebo \_CZ\_new.txt \\ \hline
		discard\_changes          & Zahodí změny a obnoví poslední načtený text \\ \hline
		on\_text\_changed         & Sleduje změny textu a nastaví příznak unsaved\_changes \\ \hline
		toggle\_original          & Přepínač mezi original\_view a exit\_original\_view \\ \hline
		original\_view            & Načte originální verzi textu aktuální stránky \\ \hline
		exit\_original\_view      & Vrátí se k upravené verzi textu aktuální stránky \\ \hline
		toggle\_translate         & Přepínač mezi start\_translate a stop\_translate \\ \hline
		start\_translate          & Přepne editor na český překlad aktuální stránky \\ \hline
		stop\_translate           & Vrátí editor na německý originál aktuální stránky \\ \hline
	\end{tabularx}
	\label{tab:texteditormanager}
\end{table}

\subsubsection{app/page\_controller.py}

Spodní navigační lišta pro přepínání stránek. Dědí z \texttt{QWidget},
emituje signál \texttt{pageChanged}.



\begin{table}[H]
	\centering
	\caption{Metody třídy PageController}
	\begin{tabularx}{\textwidth}{|l|X|}
		\hline
		\textbf{Metoda} & \textbf{Popis} \\ \hline
		\_\_init\_\_           & Vytvoří layout s tlačítky šipek a polem pro číslo stránky \\ \hline
		set\_max\_page         & Nastaví maximální číslo stránky a aktualizuje tlačítka \\ \hline
		set\_page              & Nastaví aktuální stránku a aktualizuje UI \\ \hline
		increment              & Přejde na další stránku \\ \hline
		decrement              & Přejde na předchozí stránku \\ \hline
		page\_entered          & Přečte číslo ze vstupního pole a přejde na zadanou stránku \\ \hline
		update\_buttons\_state & Deaktivuje šipku na hranicích (první / poslední stránka) \\ \hline
	\end{tabularx}
	\label{tab:pagecontroller}
\end{table}

\clearpage

\subsection{Contents Window}

\subsubsection{app/contents.py}

Definice a logika okna Contents (viz sekci~\ref{sec:contents}).

\textbf{\texttt{Ui\_Contents}} — sestavuje layout okna a propojuje komponenty.

\begin{table}[H]
	\centering
	\caption{Metody třídy Ui\_Contents}
	\begin{tabularx}{\textwidth}{|l|X|}
		\hline
		\textbf{Metoda} & \textbf{Popis} \\ \hline
		setupUi            & Vytvoří layout, načte data a naváže vyhledávání \\ \hline
		add\_year\_buttons & Naplní levý panel tlačítky ročníků \\ \hline
		show\_chapters     & Zobrazí kapitoly zvoleného ročníku \\ \hline
		chapter\_clicked   & Načte stránku a případně zvýrazní hledaný text \\ \hline
		retranslateUi      & Nastaví titulek okna \\ \hline
	\end{tabularx}
	\label{tab:uicontents}
\end{table}

\textbf{\texttt{openContentsWindow(main\_window)}} — vytvoří a zobrazí okno Contents
jako samostatný \texttt{QWidget}.

\subsubsection{app/contents\_loader.py}

Načítání obsahu z Excel souboru \texttt{Obsah\_Inhalt\_1891-1941.xls}.

\textbf{\texttt{load\_contents()}} — načte soubor pomocí \texttt{pandas}, projde řádky
a vrátí slovník \texttt{\{ rok: [\{ page, title \}] \}}.

\subsubsection{app/contents\_search.py}

Logika vyhledávání v okně Contents.

\begin{table}[H]
	\centering
	\caption{Funkce a třídy modulu contents\_search.py}
	\begin{tabularx}{\textwidth}{|l|X|}
		\hline
		\textbf{Funkce / třída} & \textbf{Popis} \\ \hline
		\_filter\_contents & Filtruje data podle vyhledávacího dotazu (case-insensitive) \\ \hline
		\_clear\_left      & Odstraní všechna tlačítka ročníků \\ \hline
		\_clear\_right     & Odstraní všechna tlačítka kapitol \\ \hline
		\_restore\_full    & Obnoví plný seznam ročníků po vymazání vyhledávání \\ \hline
		\_apply\_filter    & Aplikuje filtr a dávkově renderuje výsledky s HTML zvýrazněním \\ \hline
		LineEditFilter     & QObject event filter — spustí vyhledávání při zaostření pole \\ \hline
		bind\_search       & Naváže QTimer s 500 ms zpožděním na změny textu v lineEdit \\ \hline
	\end{tabularx}
	\label{tab:contents-search}
\end{table}

\clearpage
\subsection{Settings Window}

\subsubsection{app/settings.py}

Hlavní okno nastavení (viz sekci~\ref{sec:settings}). Emituje signál
\texttt{config\_saved} po uložení.

\textbf{\texttt{SettingsWindow}}

\begin{table}[H]
	\centering
	\caption{Metody třídy SettingsWindow}
	\begin{tabularx}{\textwidth}{|l|X|}
		\hline
		\textbf{Metoda} & \textbf{Popis} \\ \hline
		\_\_init\_\_          & Sestaví layout s panely, FontSettingsWidget a tlačítky OK / Cancel \\ \hline
		\_load\_file\_config  & Načte konfiguraci z conf.json pro porovnání při zavření \\ \hline
		\_on\_ok              & Uloží konfiguraci, emituje config\_saved a zavře okno \\ \hline
		closeEvent            & Zkontroluje neuložené změny a nabídne dialog Uložit / Zahodit / Zrušit \\ \hline
	\end{tabularx}
	\label{tab:settingswindow}
\end{table}

\subsubsection{app/font\_settings.py}

Widget nastavení písma vložený do \texttt{SettingsWindow}.

\textbf{\texttt{FontSettingsWidget}}

\begin{table}[H]
	\centering
	\caption{Metody třídy FontSettingsWidget}
	\begin{tabularx}{\textwidth}{|l|X|}
		\hline
		\textbf{Metoda} & \textbf{Popis} \\ \hline
		\_\_init\_\_            & Sestaví grid layout se seznamem písem, checkboxy a náhledem Sample \\ \hline
		\_apply\_config\_to\_ui & Nastaví výběry v seznamech a checkboxech podle konfigurace \\ \hline
		update\_sample          & Aktualizuje font náhledu a uloží volby do self.config \\ \hline
		save\_config            & Zapíše self.config do conf.json \\ \hline
		load\_config            & Načte konfiguraci z conf.json do self.config \\ \hline
	\end{tabularx}
	\label{tab:fontsettingswidget}
\end{table}
		
		
\chapter{Závěr}

V rámci semestrálního projektu se podařilo navrhnout a implementovat desktopovou aplikaci zaměřenou na práci s~digitalizovanými historickými dokumenty. Vytvořený nástroj poskytuje uživatelsky přívětivé prostředí, které efektivně propojuje vizuální podobu naskenovaných stránek s~jejich textovými přepisy. \\[6pt]

Aplikace úspěšně integruje klíčové funkce, jako je \textbf{paralelní zobrazení obrazu a textu}, \textbf{podpora bilingvního obsahu} (český a německý jazyk) a \textbf{strukturovaná navigace} v~dokumentu. Implementovaný textový editor umožňuje nejen pohodlné prohlížení, ale i \textbf{přímou editaci a opravu přepisů}, což je zásadní pro práci s~historickými prameny. Součástí aplikace je také \textbf{vyhledávání} a přehled obsahu, které výrazně usnadňují orientaci v~rozsáhlých dokumentech. \\[6pt]

Výsledná aplikace představuje funkční řešení odpovídající specifickým požadavkům oblasti digitalizace a archivnictví. Přínosem je zejména \textbf{propojení originálního zdroje s~jeho digitální verzí}, což zefektivňuje proces zpracování historických textů a zvyšuje jejich dostupnost pro další využití. \\[6pt]

Do budoucna lze aplikaci dále rozšířit o~\textbf{automatické rozpoznávání textu (OCR)} s~využitím metod umělé inteligence, které by umožnilo generování prvotních přepisů přímo z~naskenovaných obrazů. Dalším směrem vývoje je implementace \textbf{cloudového ukládání a sdílení dat}, což by umožnilo spolupráci více uživatelů na jednom projektu. Uvažovat lze také o~rozšíření podpory dalších jazyků nebo o~pokročilé nástroje pro analýzu historických textů.
	
\chapter{Bibliografie}

Při vývoji byly využity následující podpůrné zdroje:

\begin{itemize}
	\item \textbf{Nástroje AI (ChatGPT, Claude):} využity pro dílčí konzultace, optimalizaci a asistenci při ladění kódu.
	\item \textbf{Externí knihovny:} komponenta pro zobrazování dokumentů využívá jako technický základ modul \texttt{QImageViewer.py} od autora acbetter (\url{https://gist.github.com/acbetter/32c575803ec361c3e82064e60db4e3e0}).
\end{itemize}


		
	\end{document}
